% BEGIN THERMAL_R03_SYNC_13B
% Consolidated formalization of results already presented in the body.
\subsection{Structural definitions and quantitative results}
\label{sec:ii-conclusiones-cuantitativas}

Definition through structure preserves the object prior to the scalar
reading. For the action constants, that object is an oriented section
with a multiplicative norm over the local quotient and a regional character. For Barbero--Immirzi,
it is the evaluation of an oriented incidence. For mixing, it is a transport
between coordinates and amplitudes. For Boltzmann, it is a map between
lines of quantities. For gravitation, it is a reciprocal action section
whose magnitude follows from the inscribed length and radial reading
specified in R1--R8. These definitions
determine different operations within the same genealogy.

For each quantity, the following table fixes the object that defines it,
the operation that determines its reading, and the physical role developed
in the body. The reconstruction and transport proofs establish
how that information is preserved when the result is reused. In
particular, \ref{prop:registro-36-residuos} characterizes the representation
of the record, \ref{mem:prop-schur} preserves observation under reduction,
and \ref{cor:ii-witt-composed-reader} preserves the angular coordinates
in the incidence realization. These three operations have their own domains
and distinct demonstrative roles within the same sequence.

\begingroup
\setlength{\LTpre}{0pt}
\setlength{\LTpost}{2pt}
\begin{longtable}{@{}p{.13\linewidth}p{.415\linewidth}p{.395\linewidth}@{}}
\caption{Structural object, produced reading, and physical meaning of
the article's results.}\label{tab:ii-conclusiones-estructura}\\
\toprule
Quantity & Structural definition and result & Physical reading and internal development\\
\midrule
\endfirsthead
\toprule
Quantity & Structural definition and result & Physical reading and internal development\\
\midrule
\endhead
\bottomrule
\endfoot
Planck & Section of the action line:
$\hbar_{\rm pre}=H_5\,10^{-34}\mathcal U_S$,
$\hbar_{\rm ret}=\eta_{\rm ret}\,10^{-34}\mathcal U_S$.
The local quotient has sixteen sheets, and
$\operatorname{ord}_{10}(\varphi\alpha^{16})=-34$.
& The section relates frequency and energy; the complete turn determines
$h_a=2\pi\hbar_a$. Development in \ref{sec:action}; coordinates and comparison
in \ref{tab:ii-planck-comparacion}.\\[5pt]
Barbero--Immirzi & Oriented evaluation of the channels and incidence:
$\gamma_{\HMT}=\pi AC^*/180+12\Delta S_{90}-\Delta S_{120}$.
The printed evaluation is $0.2740076726944592747\ldots$.
& The coefficient enters elementary area resolution and
the Cartan--Holst realization. The definitions in
\ref{sec:barbero} and \ref{sec:area} preserve these two uses;
\eqref{eq:holst} specifies the second.\\[5pt]
CKM mixing & Image of the angular vector under the three sectorial
maps; phase $\delta_{\deg}=9A_{\deg}$.
The angles and amplitudes are evaluated after the reader is constructed.
& The Jarlskog quartet expresses phase information invariant
under diagonal basis changes. The rules, their evaluation, and the
unitary realization are developed in
\ref{sec:ii-propagacion-angular}.\\[5pt]
Boltzmann & Positive transduction
$k_B:\mathcal L_\Theta\to\mathcal L_E$ with
$k_B\Theta_{{\rm clk},a}\log3=h_a/T_{108}$.
The state and measure of the continuations determine the information
content on which it acts.
& $S=k_Bs$ dimensionally realizes the entropy of the reduced state.
Energy and temperature bases and their transformations are preserved
in \ref{sec:ii-boltzmann}; the additive characterization and instrument
are proved in \ref{sec:ii-aditiva-instrumento}; the comparison is in
\ref{tab:ii-kb-g-cartas}.\\[5pt]
Marked thermal realization & Fixed reference $\eta$, with multiplicities
$(d_0,d_1,d_2)=(1,380,649)$ at degrees $23,26,29$ and trace
$b_*=1.380649\times10^{-23}$.
The readings $\mathcal S_P=b_*s$ and $\mathcal E_P=\operatorname{Tr}(\rho H)$
determine $d\mathcal S_P/d\mathcal E_P=b_*\beta$ and
$\Theta_P=(b_*\beta)^{-1}$.
& With fixed reference, transported positive bases and a nonscalar
Hamiltonian, \ref{thm:ii-ter27-transducer} determines the transducer.
Corollary~\ref{cor:ii-ter27-variational} proves the unique minimum.
Composition preserves $q^N/Z_N(q)$, the energy origin and the
clock--horizon transports of \eqref{eq:ii-ter27-intertwiner}.\\[5pt]
Gravitation & Reciprocal section
$G_a=c_{\rm int}^{3}L^2/\hbar_a$,
with $L=(5759/23040)\alpha_{\HMT}^{16}L_*$.
The clock $t_0=\pi_{\HMT}L/(54c_{\rm int})$ preserves the circular
expression $G_a=\vartheta_{108}^{\circ}c_{\rm int}^{5}t_0^2/\hbar_a$,
with $\vartheta_{108}^{\circ}=(54/\pi_{\HMT})^2$.
& The product $G_a\hbar_a$ preserves areal normalization.
The longitudinal and radial construction, its rules, and uniqueness
of the coefficient are proved in \ref{g26:sec:rigidez-radial-en};
the circular realization is retained in \ref{sec:gravity} and thermal confluence in
\ref{sec:confluencia}.\\
\end{longtable}
\endgroup

Quantitative comparison preserves the particular meaning of each
reference. For Planck, Table~\ref{tab:ii-planck-comparacion} compares
dimensional coordinates with an SI defining constant; the relative
residuals of the pre-return and post-return orientations are,
respectively, approximately $-1.82\times10^{-15}$ and
$-6.93\times10^{-10}$. For Barbero--Immirzi,
Table~\ref{tab:ii-barbero-comparacion} compares two constructions:
the internal evaluation and the root of the Ghosh--Mitra equation, whose
approximate value is $0.2740668582328300$. The relative difference is
approximately $-2.16\times10^{-4}$; the object compared is a
theoretical coefficient.

For CKM, Table~\ref{tab:ii-ckm-comparacion} brings together the sines of the
three mixing angles, the phase in radians, and the Jarlskog invariant.
The angles and the matrix of moduli are made explicit in the adjacent
development.
The normalized marginal deviations defined there satisfy
$|z|<0.029$ relative to the PDG rows used. Each residual uses
the indicated marginal uncertainty. Its reading is individual;
a joint evaluation additionally requires the covariance matrix
of those observables.
For Boltzmann and gravitation, Table~\ref{tab:ii-kb-g-cartas}
presents the transducer, longitudinal and radial composition, dimensional charts,
and external reference separately. This organization makes it possible
to identify exactly which quantity each construction produces and
which quantity each comparison evaluates.

Table~\ref{tab:ii-conclusiones-contraste} brings together the results of
these comparisons alongside the preceding structural definitions.
\begin{table}[!htbp]
\centering\small
\setlength{\tabcolsep}{4pt}
\renewcommand{\arraystretch}{1.12}
\begin{tabularx}{\linewidth}{@{}>{\raggedright\arraybackslash}p{.16\linewidth}>{\raggedright\arraybackslash}X>{\raggedright\arraybackslash}X@{}}
\toprule
Quantity & Result in the specified realization & Reference and comparison\\
\midrule
Planck, before return
& \(\hbar_{\rm pre}=1.0545718176461544696\ldots\)
  \(\times10^{-34}\,\mathrm{J\,s}\).
& SI: \(\hbar_{\rm SI}=h_{\rm SI}/(2\pi)\), exact.
  Its coordinate is \(1.0545718176461563912\ldots\)
  \(\times10^{-34}\,\mathrm{J\,s}\).
  Relative residual \(-1.822221\times10^{-15}\).\\[4pt]
Planck, after return
& \(\hbar_{\rm ret}=1.0545718169150390611\ldots\)
  \(\times10^{-34}\,\mathrm{J\,s}\).
& Same SI reference.
  Relative residual \(-6.932836\times10^{-10}\).
  For \(h_a=2\pi\hbar_a\), the residuals are identical.\\[4pt]
Barbero--Immirzi
& \(\gamma_{\HMT}=0.2740076726944592747\ldots\).
& Ghosh--Mitra: approximately \(0.2740668582328300\).
  Relative difference \(-2.159529202\times10^{-4}\)
  between theoretical constructions.\\[4pt]
CKM mixing
& \(\delta_{\rm CP}=65.676173123554^\circ\),
  \(J=3.1189723326632974\times10^{-5}\).
  The three angles in \eqref{eq:ii-ckm-decimales} are
  \(\theta_{12}=13.003313589509^\circ\),
  \(\theta_{23}=2.398056157332^\circ\), and
  \(\theta_{13}=0.213977382244^\circ\).
& PDG 2025: all five comparisons of
  \(s_{12},s_{23},s_{13},\delta,J\) satisfy \(|z_x|<0.029\)
  with the marginal uncertainties in
  Table~\ref{tab:ii-ckm-comparacion}.\\[4pt]
Boltzmann
& \(B_*=1.380649\times10^{-23}\,u_E/u_\Theta\),
  with the bases and thermometric rule of
  \eqref{eq:ii-b-evaluacion-comparada}.
& SI: \(1.380649\times10^{-23}\,\mathrm{J/K}\), exact.
  The numerals coincide under the declared dimensional transport.\\[4pt]
Length \(L\)
& \(L=1.6162549826251263301\ldots\)
  \(\times10^{-35}\,\mathrm m\), for \(L_*=1\,\mathrm m\).
& Longitudinal and radial composition:
  \(G_{\rm ret}=c_{\rm int}^3L^2/\hbar_{\rm ret}\),
  with the same velocity and action bases.\\[4pt]
Gravitation
& \(G_{\rm ret}=6.6742996594140227\ldots\)
  \(\times10^{-11}\,\mathrm{m^3\,kg^{-1}\,s^{-2}}\).
& CODATA 2022: \(6.67430(15)\times10^{-11}\) in the same units.
  Difference of approximately \(-0.00227057\) standard uncertainties.\\
\bottomrule
\end{tabularx}
\caption{Summary of the quantitative comparison. Evaluations, bases,
and references are those in Tables~\ref{tab:ii-planck-comparacion}--\ref{tab:ii-kb-g-cartas}.
The relative residuals of the action constants, the theoretical comparison
of Barbero--Immirzi, and the normalized residuals of CKM and gravitation
retain their respective meanings.}
\label{tab:ii-conclusiones-contraste}
\par\medskip
\centering\small
\setlength{\tabcolsep}{4pt}
\renewcommand{\arraystretch}{1.18}
\begin{tabular}{@{}lrrr@{}}
\toprule
Observable & HMT evaluation & PDG 2025 & \(z_x\)\\
\midrule
\(s_{12}\) & \(0.225007404757\) & \(0.22501\pm0.00068\) & \(-0.003817\)\\
\(s_{23}\) & \(0.041841757010\) & \(0.04183^{+0.00079}_{-0.00069}\) & \(0.014882\)\\
\(s_{13}\) & \(0.003734601164\) & \(0.003732^{+0.000090}_{-0.000085}\) & \(0.028902\)\\
\(\delta\) (rad) & \(1.146265461116\) & \(1.147\pm0.026\) & \(-0.028251\)\\
\(10^5J\) & \(3.118972333\) & \(3.12^{+0.13}_{-0.12}\) & \(-0.008564\)\\
\bottomrule
\end{tabular}
\caption{CKM values and marginal comparison collected in the conclusions.
The three angles in Table~\ref{tab:ii-conclusiones-contraste} are compared
through \(s_{ij}=\sin\theta_{ij}\), together with the phase in radians and
the invariant \(J\). The values, uncertainties, and residuals from
Table~\ref{tab:ii-ckm-comparacion} are retained, corresponding to the
PDG 2025 three-generation fit. The last row expresses both the value
and its uncertainty in units of \(10^{-5}\). The residuals are marginal
comparisons of correlated observables.}
\label{tab:ii-conclusiones-ckm}
\end{table}


\subsection{Conservation of the areal scale and recovery of boundary information}

The relation between information and gravitation is concentrated in two
complementary maps. The first preserves the areal scale through
$G_a\hbar_a$ and organizes its occupations through
$\gamma_{\HMT}a_0$. The second preserves sectorial resolution through
the joint observation of area and weighted incidence. Inversion
\eqref{eq:recover} recovers the occupations from area and the
weighted observable; equations
\eqref{eq:ii-respuesta-entropica} and \eqref{eq:ii-respuesta-areal}
link their variations to fluctuations of the state.

This gives the article's central physical reading: the same
configuration admits a geometric observation and an information
observation related by a recoverable structure. An orientation
change can preserve the entropy of the reduced state and
transform area, as established by
\eqref{eq:ii-inversion-entropia} and \eqref{eq:ii-area-complementaria}.
Information conservation concerns the state and its transport;
entanglement entropy additionally concerns the bipartition and
reduction. The horizon identity in
\ref{sec:confluencia} uses these roles in the specified
geometric and thermal realization.

The memory identities also allow the scope of variable
reduction to be specified. Equality of observations in
\eqref{mem:schur-lector} preserves contributions of the eliminated sector
through the transformed effective operator, source, and reader.
The entropy of the reduced state, by contrast, is evaluated with the bipartition
and measure specified in \ref{sec:ii-ecuacion-informacion}.
The two constructions explain different roles of memory:
the first preserves a reading of transport; the second quantifies
the information accessible to the subsystem. Their articulation keeps
explicit which structure is transported and which reduction defines the
physical observable.

The explanatory advance is therefore expressed through a precise order:
generated structure, reading map, resulting value, and physical
role. Equality of two values is studied together with equality
or difference of their structures; the recovery maps determine
which of those differences retains observable meaning.
% END THERMAL_R03_SYNC_13B
